\documentclass[10pt]{article}

\usepackage[utf8]{inputenc}

\usepackage[T1]{fontenc}

\usepackage{amsmath}

\usepackage{amsfonts}

\usepackage{amssymb}

\usepackage[version=4]{mhchem}

\usepackage{stmaryrd}

\usepackage{graphicx}

\usepackage[export]{adjustbox}

\graphicspath{ {./images/} }



\DeclareUnicodeCharacter{0131}{$\imath$}



\begin{document}

семирегулярным) $T_{0}$-пространством, а верхний $\hat{S}$ и нижний $\stackrel{S}{S}$ пределы суть, $T_{1}$-пространства. Если $S-$ конечный ГІ. с., то $\bar{S}, \hat{S}$ и $\breve{S}$ - бикомпактные пространства.



В основе всей теории аппроксимации топологич. пространств полиәдрами, вернее симплициальными комплексами, лежит введенное П. С. Александровым (см. [1]) понятие нерва системы множеств. Н е р в о м данной системы $\alpha$ множеств наз. симплициальный комплекс $N_{\alpha}$, вершины к-рого взаимно однозначно соответствуют элементам системы $\alpha$ таким образом, что каждое множество вершин определяет симплекс комплекса $N_{\alpha}$ тогда и только тогда, когда соответствующие этим вершинам множества системы $\alpha$ имеют непустое пересечение.



Удобнее рассматривать т. н. канонич. покрытия пространства $X$. Локально конечное (конечное) покрытие $\alpha$ пространства $X$ наз. к а н о н и ч е с к и м, если его элементы - канонические множества (замкнутые) (в другой терминологии - регулярные замкнутые) с дизъюнктными открытыми ядрами. Если из двух канонич. покрытий пространства $X$ покрытие $\alpha^{\prime}$ следует за покрытием $\alpha$, т. е. $\alpha^{\prime}$ вписано в $\alpha$ (в этом случае $\left.\alpha^{\prime}>\alpha\right)$, то определено естественное симплициальное отображение $\pi_{\alpha}^{\alpha^{\prime}}$ (проекция) нєрва $N_{\alpha}$, на нерв $N_{\alpha}$, к-рое возникает, если каждому элементу $A^{\alpha^{\prime}} \in \alpha^{\prime}$ покрытия $\alpha^{\prime}$ поставить в соответствие тот единственный элемент $A^{\alpha}$ покрытия $\alpha$, для к-рого $A^{\alpha} \supset A^{\alpha^{\prime}}$. Пусть $\mathfrak{A}(X)$ (соответственно $\mathfrak{A}_{0}(X)$ ) обозначает совокупность всех локально конечных (конечных) канонич. покрытий пространства $X$. Для каждого $\alpha \in \mathfrak{Y}(X)$ (соответственно $\left.\alpha \in \mathfrak{U}_{0}(X)\right)$ обозначена вся совокупность всех локально конечных (конечных) канонич. покрытий пространства $X$. Для каждого $\alpha \in \mathfrak{Y}(X)$ (соответственно $\alpha \in \mathfrak{A}_{0}(X)$ ) рассматривается комплекс $N \alpha$, являющийся нервом покрытия $\alpha$. Если $\alpha^{\prime}>\alpha$, то определено симплициальное отображение $\pi_{\alpha}^{\alpha^{\prime}}: N_{\alpha^{\prime}} \rightarrow N_{\alpha}$. Полученный таким образом П. с. $S=\left\{N_{\alpha}, \pi_{\alpha}^{\alpha^{\prime}}\right\}$ наз. по л н ы м (соответственно к о н е ч н ы м) П. с. тонологич. пространства $X$. П. С. Александров [2] еще в 1928 доказал, что каждый метрический ( $n$-мерный) комıакт является верхним пределом ( $n$-мерного) конечного П. с. над счетным множеством индексов. А. Г. Курош в 1934 доказал, что каждый бикомпакт есть верхний предел своего конечного П. с. В 1961 В. И. Пономарев доказал, что каждый паракомпакт есть верхний предел своего полного П. с., то есть спектра, построенного над множеством $\mathfrak{U}(X)$ всех локально конечных канонич. покрытий пространства $X$. В. И. Пономарев ввел понятие р а с с л а б л е н и я симплициального комплекса $K$, понимая под этим всякийі замкнутый подкомплекс $K^{\prime} \subset K$, содержаций все вершины комплекса $K$. Нульмерный комплекс, состоящий из всех вершин комплекса $K$, наз. его п о л н ы м р а с с л а б л ен и е м (или о с т о в м м. Заменяя все комілексы данного П. с. их (полным) расслаблением и сохраняя при этом проекции, получают (полное) расслабление спектра. Исследование неприводимых совершенных отображений паракомпактов сводится к исследованию расслаблений их полных П. с. При этом предел полного расслабления полного $\Pi$. с. наракомпакта $X$ есть т. н. абсолют $\dot{X}$ паракомпакта $X$, а предел полного расслабления конечного П. с. любого регулярного пространства - бикомпактное расширение Стоуна - Чеха $\beta \dot{X}$ абсолюта $\dot{X}$ этого регулярного пространства. Всякий конечный абстрактный П. с. эквивалентен спектру над нек-рым направленным измельчающимся множеством конечных канонических иокрытий нек-рого полурегулярного бикомпактного $T_{0}$-пространства, т. е. получается из этого спектра посредством конечғого числа следующих операций: 1) замена спектра изоморфным ему спектром, 2) замена спектра его конфинальной частью, 3) замена спектра спектром, содержащим данный в качестве конфинальной части (т е о ре ма 3 а й це в а).



Понятия нерва и П. с. доставили средства для редукции свойств общих пространств, и прежде всего паракомпактов, бикомпактов и метрич. комиактов, к свойствам комплексов и их симплициальных отображений. Это позволило определить и изучать гомологические и когомологические инварианты общих пространств, а не только полиэдров (см. АлександроваЧеха гомологии и когомологии, Спектральные гомологии). Все это привело к синтезу геометрических и теоретикомножественных идей в топологии.



Лum.: [1] А л е к с а н д p o в П. C., "C.r. Acad. sci.», 1927, t. 184, p. 317-20; [2] e r o ж e, "Ann. Math.", 1929, v, 30, p. c. $5-57 ;$ [4] е го ж е, Введение в теорию множеств и общую с. $5-57 ;[4]$ е го ж е, Введение в теорию множеств и общую

топологию, м., 1977 ; [5] А л е к с а н д р о в П. С., П о н ом a p ё в В. И., в кн.: General topology and its relations to modern analysis and algebra, v. 2, Prague, 1967, p. 25-30; [6] А л е к с а н д р о в П. С., Ф е д о р ч ук В. В., "Успехи ма-

тем. наук", 1978, т. 33, в. 3, с. $3-48 ;[7]$ П о н м а р в в. И., "Матем. сборник», 1963 , т. 60 , № 1 , с. $89-119 ;$ [8] 3 а й це в В. И., "Труды Моск. матем. об-ва", 1972 , т. 27, с. $129-93$.

В. И. Пономарев.



ПРОЕКЦИЯ - термин, связанный с операцией п р ое к т и ровани я (п р о ци ро вани я), к-рую можно определить следующим образом (см. рис.): выбирают произвольную точку $S$ пространства в качестве це нт р a п р о е к т ир о в а ни я и илоскость П', не проходящую через точку $S$, в качестве п ло с к ост и п рое кци й. Чтобы спроектиро-



вать точку $A$ (п р о о 6 р а з) пространства на плоскость $\Pi^{\prime}$ через центр проекций $S$, проводят прямую $S A$ до ее пересечения в точке $A^{\prime}$ с плоскостью П'. Точку $A^{\prime}($ о б р а в) наз. п р о е к ц и е й точки $A$; проекцией фигуры $F$ наз. совокупность П. всех ее точек. Описанная П. наз. ц е н т р а л ь н о й (или к о в и ч е с к о ї). П. с бесконечно удаленным центром проектирования наз. п а р а л л е л ь но й (или ц и л и нд р й е с ко й). Если плоскость П. расположена перпендикулярно к направлению проектирования, то II. наз. о р т о г о а л ь н о й (или п р ям о у го л ь н о й).



Параллельные П. широко используются в начертательной геометрии для получения различных видов изображений (см., напр., Аксонометрия, Перспектива). Имеются специальные виды П1. на плоскость, сферу и др. поверхности (см., напр., Картограбическая проекция, Стереографическая проекция). А. Б. Иванов.



\includegraphics[max width=\textwidth, center]{2023_07_08_3b2bebdcc5d64b680f02g-1}

к а т е г о р и и - понятие, описывающее на языке морфизмов конструкцию декартова произведения. Пусть $A_{i}, \quad i \in I,-$ индексированное семейство объектов категории $\Re$. Объект $P \in \mathrm{Ob} \mathfrak{R}$ (вместе с морфизмами $\left.\pi_{i}: P \rightarrow A_{i}, i \in I\right)$ наз. произведением семейства объектов $A_{i}, i \in I$, если для всякого семейства морфизмов $\alpha_{i}: X \rightarrow A_{i}, i \in I$, супествует такой единственный морфизм $\alpha: X \rightarrow P$, что $\alpha \pi_{i}=\alpha_{i}, i \in I$. Морфизмы $\pi_{i}$ наз. п роек ция ми п ро и з вед ения; П. обозначается $\prod_{i \in I}^{\times} A_{i}\left(\pi_{i}\right)$, или $\prod_{i \in I} A_{i}$, или $A_{1} \times \ldots \times A_{n}$ в случае $I=\{1, \ldots, n\}$. Морфизм $\alpha$, входящий в определение П., иногда обозначается $\prod_{i \in I} \alpha_{i}$ или $(\times)_{i \in I} \alpha_{i}$. П. семейства $A_{i}, i \in I$, определено однозначно с точностью до изоморфизма; оно ассоциативно и ком-





\end{document}